\subsection{Potência em Wi-Fi}
O dBm expressa uma potencia absoluta, com 1~mW como referencia; e o numero que o celular ou o adaptador Wi-Fi reportam como intensidade de sinal. A sensibilidade do receptor e a menor potencia que ele ainda consegue decodificar, e a folga do enlace e a diferenca entre a potencia recebida e essa sensibilidade. Quatro mecanismos degradam essa potencia ao longo do trajeto: a atenuacao com a distancia, prevista por modelos de espaco livre; a absorção por paredes, mobiliario e pessoas; as reflexoes, que chegam fora de fase e podem cancelar parte do sinal (multipercurso) e o piso de ruido elevado por outros transmissores .O ganho de uma antena, medido em dBi, e a razao entre a intensidade irradiada em uma direcao e a que uma antena isotropica hipoteticairradiaria com a mesma potencia total. Esse ganho so vale dentro do lobulo principal a regiao angular em que a potencia nao cai mais de 3~dB abaixo do maximo da propria antena e, quanto maior o ganho, mais estreito e esse lobulo, tornando o enlacemais sensivel ao alinhamento (TANENBAUM; FEAMSTER; WETHERALL, 2021). Antenas omnidirecionais distribuem o ganho de forma parecida em todas as
direcoes horizontais; antenas direcionais concentram esse mesmo ganhoem poucas direcoes, ao custo das demais.

Um modelo de propagacao de espaco livre preve um decaimento suave e regular com a distancia, mas serve apenas como referencia de projeto:dentro de um predio, paredes, mobiliario, pessoas e ultipercursoacrescentam perda e variabilidade que o modelo nao captura, e o decaimento medido tende a ser maior e irregular.Por isso, nenhuma escolha de antena ou de ponto de acesso corrige, por si so, um ambiente desfavoravel: apenas a medicao descreve o ambiente
de fato.
\subsection{Mapa de Calor}
\label{subsec:heatMap}
Um mapa de calor de cobertura representa, sobre a planta do ambiente, a potencia recebida medida em uma malha de pontos. Diferentemente do padrao de radiacao do catalogo, liso e simetrico, medido sem nada em volta, o mapa de calor captura o padrao instalado: recortado por paredes, com reforco e cancelamento causados por reflexoes em superficies metalicas proximas. A malha de pontos deve ter espacamento regular e ser adensada onde o sinal muda rapidamente, proximo a paredes, portas e ao proprio ponto de acesso e nao apenas onde e comodo medir. Como o sinal recebido em um mesmo ponto oscila varios decibeis ao longo do tempo por efeito do multipercurso, uma unica leitura descreve o instante, nao o ponto; por isso e recomendado coletar multiplas amostras por ponto e registrar a mediana, que resiste melhor a uma leitura pontualmente anomala do que a media. Um mapa de calor isolado apenas descreve o ambiente. Dois mapas comparaveis construidos com a mesma malha, o mesmo aparelho e as mesmas condicoes de uso do ambiente, antes e depois de uma mudanca, permitem testar uma hipotese, como a de que e posicionar o ponto de acesso melhora a cobertura. E esse o metodo empregado neste trabalho.